% Información para el motor (arara). Compilar en LuaLaTeX o XeLaTeX sólamente.

%arara: lualatex: {
%arara: --> interaction: nonstopmode,
%arara: --> options: ['-file-line-error'],
%arara: --> shell: yes,
%arara: --> synctex: yes,
%arara: --> }
%arara: makeglossaries
%% arara: makeindex	

\documentclass[letterpaper,twoside]{article}

%%%%%%% VARIABLES ESTÁTICAS %%%%%%%
% Aqui se definen las variables para diferentes procedimientos.
% Verificar que las unidades coincidan si se deciden cambiar.

\newcommand{\defglo}[1]{\item[{\bfseries \glsname{#1}}] \glsdesc{#1}} % Define un comando para dar un item en un entorno de descripción para una palabra definida en el glosario
% \includeonly{src/PSA-2.19.tex,src/PSA-Gloss.tex}
\newcommand{\TipoID}{TEMP} % Definirlo como: PRO, PROG, ToC, G, L, CE, CI, I, etc.
\newcommand{\Prog}{TEMP}
\newcommand{\ProgTitulo}{TEMP}

\newcommand{\ie}{\emph{i.e.}}
\newcommand{\eg}{\emph{e.g.}}

%%%%%%%%%%%%%%%% Formularios %%%%%%%%%%%%%%%%
\newcommand{\IdFormAACC}{F-OP-40}

\newcommand{\BitES}{Bitácora de entradas y salidas --- Clientes (F-AD-1)}
\newcommand{\BitULl}{Bitácora de uso de llaves (F-AC-1)}
\newcommand{\BitVig}{Bitácora de vigilancia (F-AC-2)}
\newcommand{\Fpap}{Papeleta (F-AC-3)}
\newcommand{\Oent}{Orden de entrada (F-AC-4)}
\newcommand{\Osal}{Orden de salida (F-AC-5)}
\newcommand{\LayES}{Layout de carga o descarga (F-AC-6)}
\newcommand{\BitClor}{Bitácora de cloración de agua (F-AC-7)}
\newcommand{\BitTap}{Bitácora de concentración de tapete sanitario (F-AC-8)}
\newcommand{\BitUQ}{Bitácora de uso de productos químicos (F-AC-9)}
\newcommand{\BitLAd}{Bitácora de limpieza de aduana sanitaria (F-AC-10)}
\newcommand{\BitLOf}{Bitácora de limpieza de oficinas (F-AC-11)}
\newcommand{\BitLEx}{Bitácora de limpieza de exteriores (F-AC-12)}
\newcommand{\BitBPH}{Bitácora de cumplimiento de BPH (F-HyS-1)}
\newcommand{\BitCAP}{Bitácora de consumo de agua potable (F-HyS-2)}
\newcommand{\BitCE}{Bitácora de consumo de energía eléctrica (F-HyS-3)}
\newcommand{\BitIEx}{Bitácora de inspección de extintores (F-HyS-4)}
\newcommand{\BitVIR}{Bitácora de verificación de termómetros infrarrojos (F-HyS-5)}
\newcommand{\LVCMo}{Lista de verificación de condición de montacargas (F-HyS-6)}
\newcommand{\RTAC}{Registro de temperaturas con termómetro IR --- Aseguramiento de calidad (F-HyS-7)}
\newcommand{\RTMa}{Registro de temperaturas con termómetro IR --- Mantenimiento (F-HyS-8)}
\newcommand{\RAcL}{Registro de actividad de limpieza (F-HyS-9)}
\newcommand{\RAcM}{Registro de actividad de mantenimiento (F-Hys-10)}
\newcommand{\ReEq}{Registro de revisión de equipos de enfriamiento (F-MAN-1)}
\newcommand{\RCMi}{Registro de inspección de infraestructura en prevención (F-MAN-2)}
\newcommand{\Rcap}{Registro de asistencia a capacitación (F-MAN-3)}
\newcommand{\REvDQ}{Registro de eventos de derrame de químicos (F-MAN-4)}
\newcommand{\InvVyPD}{Registro de inventario de vidrio y plástico duro (F-MAN-5)}
\newcommand{\RVerSGC}{Registro de verificación del sistema de gestión de calidad (F-MAN-6)}
\newcommand{\RRCA}{Registro de recorrido de control de alérgenos (F-MC-1)}
\newcommand{\RAC}{Registro de acción correctiva o preventiva (F-OP-1)}
\newcommand{\RACQ}{Registro de acción correctiva --- Quejas (F-OP-2)}
\newcommand{\RVSCP}{Registro de verificación del servicio de control de plagas (F-OP-3)}
\newcommand{\RevD}{Registro de revisión diaria (F-SEG-1)}
\newcommand{\RCVyPD}{Registro de evento de contaminación por VyPD (F-SEG-2)}
\newcommand{\REvCA}{Registro de evento de contaminación por alérgeno (F-SEG-3)}



%%%%%%%%%%%%%%%% CÓDIGOS PSA %%%%%%%%%%%%%%%%
% \newcommand{\PSA}{PSA-3-PRO}                 % PSA-2.3 Inspección de producto durante su almacenamiento

%%%%%%%%%%%%%%%% PSA-2.3 | Inspección de producto durante su almacenamiento %%%%%%%%%%%%%%%%
\newcommand{\TiempoAndenCong}{\qty{\leq 30}{\minute}}
\newcommand{\TiempoAndenRefri}{\qty{\leq 30}{\minute}}
\newcommand{\TiempoAveriaRefri}{4}              % Tiempo máximo de permanencia de alimento refrigerado en cámara averiada (en horas) 
\newcommand{\TiempoAveriaConge}{10}             % Tiempo máximo de permanencia de alimento congelado en cámara averiada (en horas)
\newcommand{\EspacioMinimoParedProducto}{30}    % Espacio entre tarima y pared (en centimetros)
\newcommand{\EspacioMinimoTechoProducto}{30}    % Espacio entre tarima y techo (en centimetros)
\newcommand{\VecesTempManual}{dos}              % Veces que se verifica la temperatura manualmente



%%%%%%%%%%%%%%%% PSA-2.3 | Política de aprobación de proveedores %%%%%%%%%%%%%%%%
\newcommand{\NotifAuditoriaAProveedor}{2 semanas}           % Plazo máximo para notificar a proveedor sobre auditoría.
\newcommand{\PlazoResultadoAProveedor}{2 semanas}           % Plazo máximo para notificar a proveedor sobre resultados de auditoría.
\newcommand{\PlazoPlanDeAccionAuditProveedor}{2 semanas}    % Plazo maximo para recibir plan de acción de parte del proveedor tras auditoría

%%%%%%%%%%%%%%%% PSA-2.15 | Almacenamiento de registros %%%%%%%%%%%%%%%%
\newcommand{\VigenciaAlmacRegistros}{2 años}
\newcommand{\VigenciaAlmacRegistrosElec}{3 años}

%%%%%%%%%%%%%%%% GLOBAL | Resetear glosarios en cada sección %%%%%%%%%%%%%%%%
\usepackage{etoolbox}
\preto\section{\glsresetall}

\newcommand{\AC}{aseguramiento de calidad}
\newcommand{\G}{gerencia}
\newcommand{\AD}{alta dirección}
\newcommand{\OP}{operaciones}
\newcommand{\Emb}{embarques}
\newcommand{\MC}{mesa de control}	% Definición de variables estáticas
									% i.e. Tiempo máximo de cámras de refrigeración sin energía

%%%%%%%%%%%%%%%%%%%%%%%%%%% ESPECIALES %%%%%%%%%%%%%%%%%%%%%%%%%%%%%%%%%%
% \RequirePackage{expl3}    % \usepackage cannot be used before \documentclass
% \ExplSyntaxOn             % Switch on expl3 syntax
% \pdf_version_gset:n{1.7}  % Use provided expl3 function
% \ExplSyntaxOff            % Switch off expl3 syntax

%%%%%%%%%%%%%%%%%%%%%%%%%%% Tipografía %%%%%%%%%%%%%%%%%%%%%%%%%%%%%%%%%%

\usepackage[sfdefault,t,lf]{carlito}    % Alternativa OpenSource a Calibri, definir sans serif como predeterminado, numerales tabulares
% \usepackage[default,lnum]{opensans}
% \usepackage{CormorantGaramond}

% \usepackage{lucidabr}

% \usepackage{fontspec}

\fontsize{8.8pt}{12pt}\selectfont
\renewcommand{\baselinestretch}{1.2}

% \linespread{1.3}
% \DeclareUnicodeCharacter{2212}{-}

\usepackage{polyglossia}
	\setdefaultlanguage[spanishoperators=all]{spanish}  % Definir la localización del documento como es-MX (para que cuadro sea tabla, etc.)
	% \setdefaultlanguage{mexican}
	\renewcommand{\listtablename}{Índice de tablas}

\usepackage{changelog}  % Control de versiones

\usepackage{pdflscape}
\usepackage{pifont}
\newcommand{\cmark}{\ding{51}}%
\newcommand{\xmark}{\ding{55}}%

\usepackage[%
toc=true,
postpunc={;},
acronyms,
xindy]{glossaries-extra}    % Manager de glosarios, usar xindy para el indice alfabético
% \usepackage{glossary-mcols}
% \renewcommand{\glsmcols}{2}
\setglossarystyle{long3col}

	\makeglossaries{}
	% \setabbreviationstyle[acronym]{short-long}  % Estilo de acrónimos
	\setabbreviationstyle[acronym]{long-short}
	% \setglossarystyle{mcolindex}
	\loadglsentries[main]{../glosario-def.tex}  % Archvio que LuaLaTeX toma para las definiciones

	\renewcommand{\glossarypreamble}{\fontsize{10pt}{12pt}\selectfont}

	
\usepackage{imakeidx}       % Para hacer indices
	\makeindex[columns=2]   % Estilo del indice alfabético

\newcounter{note}[section]  % Nuevo contador para notas
\renewcommand{\thenote}{\thesection.\Alph{note}}    % Definir estilo de numeración de las notas
\usepackage[framemethod=TikZ]{mdframed}             % Motor gráfico para notas (TikZ)
\newenvironment{note}[1][]{%                        % Nuevo entorno de Nota (note) 
    \refstepcounter{note}
    \begin{mdframed}[%                                  
        frametitle={Nota \thenote\ #1},                 % Titulo de las notas
        skipabove=\baselineskip{} plus 2pt minus 1pt,     % Espaciado superior
        skipbelow=\baselineskip{} plus 2pt minus 1pt,     % Espaciado inferior
        linewidth=0.5pt,                                % Tamaño de linea del recuadro
        frametitlerule=true,                            % incluir titulo
        frametitlebackgroundcolor=gray!30               % Color del recuadro
    ]%
}{%
    \end{mdframed}
}

\newcounter{contact}[section]                           % Nuevo contador para contactos
\renewcommand{\thecontact}{\thesection.\Alph{contact}}  % Definir que la numeración sea Sección.Alfabético

\newenvironment{contact}[1][]{%
    \refstepcounter{contact}
    \begin{mdframed}[%
        frametitle={Contacto \thecontact\ #1},          % Titulo del contacto
        skipabove=\baselineskip{}plus 2pt minus 1pt,	% Espaciado superior
        skipbelow=\baselineskip{}plus 2pt minus 1pt,    % Espaciado inferior
        linewidth=0.5pt,                                % Tamaño de linea del recuadro
        frametitlerule=true,                            % poner titulo
        frametitlebackgroundcolor=yellow!30             % Color del recuadro
    ]%
}{%
    \end{mdframed}
}

\usepackage[multiple]{footmisc}     % Cambios en la presentación de pies de página
\usepackage{microtype}              % Ajustes de kerning y otros 
\usepackage{graphicx}               % Mejora al motor de imagenes
\usepackage{hyperref}               % Referencias cruzadas en PDF
\usepackage{cleveref}               % Referencias cruzadas en texto
%%%%%% REDEFINICIONES DE REFERENCIAS CRUZADAS %%%%%%%%%
    \crefname{note}{nota}{notas}    
	\Crefname{note}{Nota}{Notas}
	\crefname{contact}{contacto}{contactos}
	\Crefname{contact}{Contacto}{Contactos}
	\crefname{paragraph}{parrafo}{parrafos}
	\Crefname{paragraph}{Parrafo}{Parrafos}
%%%%%%%%%%%%%%%%%%%%%%%%%%%%%%%%%%%%%%%%%%%%%%%%%%%%%%%
\newcommand{\email}[1]{\href{mailto:#1}{#1}}    % Comando para hacer que se puedan clickear emails
\newcommand{\cels}[1]{\qty{#1}{\celsius}}       % Escrir celsius rapido

\usepackage{lastpage}   % Definir variable para el numero total de páginas
\usepackage{siunitx}    % Unidades internacionales y alineamiento de decimales en tablas
	\sisetup{detect-all}    % Detectar automaticamente que la tipografía es sans serif

\usepackage{color}  % definir colores nuevos

\usepackage{tabularray} % Creación de tablas
	\UseTblrLibrary{booktabs,siunitx}   % librerias parea tablas correctas, unidades internacionales
	\definecolor{Gallery}{rgb}{0.929,0.929,0.929}   % definición del color "Gallery"
	\DefTblrTemplate{contfoot-text}{normal}{\textit{Continúa en la siguiente página}}   % Continuación de tabla
	\SetTblrTemplate{contfoot-text}{normal} % Continuación de tabla normal (footer)
	\DefTblrTemplate{conthead-text}{normal}{\textit{(Continuación)}}    % Continuación de tabla
	\SetTblrTemplate{conthead-text}{normal} % Continuación de tabla (header)

    % \usepackage[margin=1cm,bmargin=2cm,tmargin=5cm,headheight=110pt]{geometry}   % Dispoción espacial del documento
% \newgeometry{includeheadfoot, heightrounded, left=1cm, right=1cm, top=1cm, bottom=4cm, headheight=6em, footskip=1.5em}    
%     \savegeometry{L1}

% \newgeometry{includeheadfoot, heightrounded, left=1cm, right=1cm, top=1cm, bottom=4.5cm, headheight=7em, footskip=1.5em, showcrop, showframe}
%     \savegeometry{L2}

\usepackage[%
includeheadfoot,
heightrounded,
left=1cm,
right=1cm,
top=1cm,
bottom=1cm,
headheight=110pt,
footskip=1.5em]{geometry}    

\usepackage[export]{adjustbox}

\usepackage{enumitem}
	\setlist[1]{labelindent=\parindent} % Indentación como parrafo para listas
	\setlist[description]{style=nextline}
	\setlist[itemize]{leftmargin=*} 	% Eliminar margen en izq
	\setlist[itemize,1]{label=---}  	% Viñeta como em-dash
	\setlist[itemize,2]{label=--}
	\setlist{noitemsep} 				% Que no se separen los item
	% [label={\theparagraph--\arabic*}] % Numerarar parrafo-numero

	\newlist{ListaSec}{enumerate}{3}
	\setlist[ListaSec]{label* = \arabic*., name* = \arabic{section}.\arabic{subsection}}


%%%%%%%%%%%%%%%% DEFINICIÓN DE CABECERA %%%%%%%%%%%%%%%%
\usepackage{fancyhdr}   % Usar fancy headers y footers
	\renewcommand{\headrulewidth}{0pt}  % Grosor de regla del header
	\fancyhead[CE,CO,LE,LO,RE,RO]{} % Resetear headers    
    \fancyfoot[CE,CO,LE,LO,RE,RO]{} % Resetear footers

%%%%%%%%%%%%%%%%%%%%%% VARIABLES DE CADA DOCUMENTO DENTRO DE UN PROGRAMA %%%%%%%%%%%%%%%%%%%%%
	\newcommand{\Codigo}{}                  % Codigo del documento
	\newcommand{\FechaPub}{}                % Fecha de publicación
	\newcommand{\Titulo}{}                  % Titulo del documento
	\newcommand{\MayorVer}{}                % Versión mayor
	\newcommand{\MenorVer}{}                % Verisón menor
	\newcommand{\Edit}{\MayorVer.\MenorVer} % Definición de la edición

    \newcommand{\Elaboro}{}                 % Nombre de quien lo elaboró
    \newcommand{\Reviso}{}                  % Nombre de quien revisó el documento
    \newcommand{\Autorizo}{}                % Nombre de quien aprobó el documento

%%%%%%%%%%%%%%%%%% Formato por defecto para informaicón documentada general %%%%%%%%%%%%%%%%%%
\fancypagestyle{formato-PS}{%
	\fancyhead{}
	\fancyfoot{}
	\fancyhead[RE]{%
		\begin{tblr}{%	
			colspec = {X[c,m]X[c,m]},    % Especificación de columnas: Anchura estática, centrado en medio
			hlines,
			vlines,
			width = 0.35\linewidth,}%
			{\textbf{Código:}\\ \Codigo} 				 & {\textbf{Edición:}\\ \MayorVer.\MenorVer} 			\\ 	% Logotipo (2 columnas), red de frios (3 columnas), código del documento
			{\textbf{Fecha de publicación:}\\ \FechaPub} & {Página~\thepage~de~\pageref{LastPage}} 		% Titulo (4 columnas), Fecha, edición
		\end{tblr}
	}%

	\fancyhead[LO]{%
		\begin{tblr}{%	
			colspec = {X[c,m]X[c,m]},   % Especificación de columnas: Anchura estática, centrado en medio
			hlines,
			vlines,
			width = 0.35\linewidth,}%
			{\textbf{Código:}\\ \Codigo} 				& {\textbf{Edición:}\\ \MayorVer.\MenorVer} 			\\	% Logotipo (2 columnas), red de frios (3 columnas), código del documento
			{\textbf{Fecha de publicación:}\\\FechaPub}	& {Página~\thepage~de~\pageref{LastPage}} 		% Titulo (4 columnas), Fecha, edición
		\end{tblr}
	}%

	\fancyfoot[C]{Página~\thepage~de~\pageref{LastPage}}   % Para que muestre Página # de #
	\fancyfoot[L]{\textbf{CONFIDENCIAL}}                    % CONFIDENCIAL
	\fancyfoot[R]{\Codigo~v.\Edit}                          % Estructura del código CODIGO v.EDICION
}%

\fancypagestyle{formato-PI}{%
	\fancyhead{}
	\fancyfoot{}
    \fancyhead[C]{%
	    \begin{tblr}{%	
            colspec = {X[c,m]X[c,m]X[c,m]X[c,m]X[c,m]X[c,m]},    % Especificación de columnas: Anchura estática, centrado en medio
            hlines,
            vlines,
			width = \linewidth,}%
			\SetCell[c=2]{} \includegraphics[height=3em, valign=c]{./RDF_Logo.png}	&	& \SetCell[c=3]{} \large{Red de fríos S.A. de C.V.}	&	&												& {\textbf{Código:}\\ \Codigo}	\\ % Logotipo (2 columnas), red de frios (3 columnas), código del documento
			\SetCell[c=4]{} {\textbf{Titulo:}\\ \Titulo} 							&	&													&	& {\textbf{Fecha de publicación:}\\ \FechaPub}	& {\textbf{Edición:}\\ \Edit}	\\ % Titulo (4 columnas), Fecha, edición
            \end{tblr}
	}
    \fancyfoot[C]{Página~\thepage~de~\pageref{LastPage}}   % Para que muestre Página # de #
    \fancyfoot[L]{\textbf{CONFIDENCIAL}}                    % CONFIDENCIAL
    \fancyfoot[R]{\Codigo~v.\Edit}                          % Estructura del código CODIGO v.EDICION
}%

%%%%%%%%%%%%%%%%%%%%%%%%%%%%%%%%%%%%%%%%%%%%%%%%%%%%%%%%%%%%%%%%%%%%%%%%%%%%%%%%%%%%%%%%%%%%%%%%%%%%%%%%%%%%%%%%%%
			
\setcounter{tocdepth}{4}            % Qué tanto mostrará la tabla de contenidos
\setcounter{secnumdepth}{6}         % Que tantas secciones numerará LaTeX (Se ajustó hasta subparrafo)

 % Define un comando para dar un item en un entorno de descripción para una palabra definida en el glosario

% \includeonly{src/BPM-1.1.tex}
\renewcommand{\Prog}{FD}

\begin{document}
\pagestyle{formato-PS}
\thispagestyle{formato-PI}
\renewcommand{\MayorVer}{1}
\renewcommand{\MenorVer}{0}
\renewcommand{\Codigo}{FD-ToC}
\renewcommand{\FechaPub}{2023--01}
%\renewcommand{\Edit}{01}
\renewcommand{\Titulo}{Programa de \textit{Food Defense} --- Índice}

\tableofcontents
\listoftables
\thispagestyle{formato-PI}
\renewcommand{\MayorVer}{2}
\renewcommand{\MenorVer}{0}
\renewcommand{\Codigo}{FD-1-PLA}
\renewcommand{\FechaPub}{2023--01}
%\renewcommand{\Edit}{2.1}
\renewcommand{\Titulo}{Programación de las reuniones de \textit{Food Defense}}
\begin{landscape}

    \vspace*{\fill}

    \section*{\Titulo} \label{PRG:FoodDefense}
    \index{Programa!Defensa Alimentaria, de!Programación}

    \begin{longtblr}[
        label = tbl:FD-PRG,
        caption = Programa de \textit{Food Defense.}
        ]{
        width = \linewidth,
        colspec = {Q[150]Q[54]Q[62]Q[58]Q[46]Q[54]Q[52]Q[44]Q[58]Q[85]Q[65]Q[81]Q[81]},
        cells = {c,m},
        row{2} = {Gallery},
        rowhead = {1},
        }
        \toprule
        Actividad                & Enero & Febrero & Marzo & Abril & Mayo & Junio  & Julio & Agosto & Septiembre & Octubre & Noviembre & Diciembre \\ \midrule
        Revisión de programa     &       &         &       &       &      & \cmark &       &        &            &         &           &           \\
        Verificación de programa &       &         &       &       &      &        &       &        &            &         &           & \cmark    \\
        \bottomrule
    \end{longtblr}

    \vspace*{\fill}

\end{landscape}	
\thispagestyle{formato-PI}
\renewcommand{\MayorVer}{2}
\renewcommand{\MenorVer}{0}
\renewcommand{\Codigo}{FD-1-MAN}
\renewcommand{\FechaPub}{2023--01}
%\renewcommand{\Edit}{2.1}
\renewcommand{\Titulo}{Manual de Defensa Alimentaria}

\section{\Titulo} \label{MAN:FoodDefense}\index{Programa!Defensa Alimentaria, de!Manual}
\subsection{Objetivo}
Establecer reglas y procedimientos, que asegure que se tornan las medidas necesarias para evitar cualquier actividad o situación voluntaria o involuntaria que pueda poner en riesgo directa o indirectamente la seguridad de los productos alimenticios que en esta empresa se manejan.

\subsection{Alcance}
\begin{itemize}
	\item Este documento detalla los lineamientos considerados por el programa de \gls{FD} de \gls{RDF}.
\end{itemize}

\subsection{Terminología y definiciones}
\begin{description}
	\item[\glsname{FD}] \glsdesc{FD} 
\end{description}

\subsection{Introducción}
El sitio debe preparar, implementar y mantener un protocolo de defensa alimentaria que se basa en una evaluación de la amenaza de la defensa alimentaria y describe los métodos, las responsabilidades y los criterios para prevenir la adulteración de los alimentos causada por actos deliberados de sabotaje. Este plan debe ser probado y revisado, como mínimo,
anualmente. El proveedor debe designar a un miembro de la alta dirección que sea responsable de la defensa alimentaria. Esta persona responsable debe asegurarse de que existen procedimientos para registrar y controlar el acceso a las áreas del sitio por parte de los empleados, contratistas y visitantes.

El plan de defensa alimentaria deberá ser probado o cuestionado para asegurar que las estrategias de defensa alimentaria son efectivas, apropiadas y documentadas. Las pruebas del plan de defensa alimentaria pueden basarse en la descripción de un escenario ficticio.

El protocolo debe identificar cómo el proveedor limita el acceso a las áreas designadas de la operación sólo a los empleados debidamente autorizados. El proveedor debe aplicar medidas para proteger los puntos de procesamiento sensibles de la contaminación intencionada. El protocolo debe explicar cómo la empresa garantiza el almacenamiento y el transporte seguros
de las materias primas, los envases, los equipos, los productos químicos peligrosos y el producto
acabado.

El protocolo debe definir cómo se van a abordar las zonas restringidas, controladas o de difícil acceso. El proveedor es libre de desarrollar medidas adecuadas para abordar áreas específicas para garantizar el control a través de una amplia variedad de soluciones.

\subsubsection{Registros}
Los siguientes son ejemplos de registros y/o documentos para ayudar a la aplicación y revisión de este tema: 

\begin{itemize}
	\item El plan de defensa alimentaria, que incluye el equipo de defensa alimentaria, la evaluación de la amenaza, las estrategias preventivas, los contactos y el plan de acción.
	\item Resultados de la revisión anual del plan de defensa alimentaria.
	\item Resultados de la prueba anual, o simulacro, al plan de defensa alimentaria.
	\item Hojas de registro, etiquetas de identificación u otros medios para controlar a los visitantes y empleados.
	\item Medidas correctivas de cualquier incumplimiento del plan
\end{itemize}

\subsection{Documentos y/o normas relacionados}
\begin{itemize}
	\item Procedimiento de Orden de entradas y Orden de salidas;
	\item Control de llaves de equipos;
	\item Procedimientos de entrada y salida de visitantes;
	\item Registro de entrada de clientes.
\end{itemize}

\subsection{Procedimiento}\index{Programa!Defensa Alimentaria, de!procedimientos}
Existe control de acceso para ingresar a oficinas y áreas administrativas de \gls{RDF} y acceso a las áreas Operativas de \gls{RDF}.
Reglamento de ingresos a \gls{RDF}
Requisitos de ingreso de entradas y salidas para visitantes y clientes.

\begin{itemize}
	\item El acceso a las áreas operativas y otros puntos cuentan con un sistema de circuito cerrado (CCTV) que trabajan las \qty{24}{\hour}.
	\item Se cuenta con vigilancia durante todo el turno laboral.
	\item El área de vigilancia controla el ingreso de proveedores y visitantes a las instalaciones previamente autorizados por el departamento correspondiente, los controles de acceso son llevados mediante registros.
	\item Cualquier persona autorizada para ingresar a las instalaciones deberá ser acompañada por personal de \gls{RDF} S.A de C.V.
	\item Registrarse en la bitácora de entrada y salida de visitantes el cual consta de los siguientes puntos:
	\begin{enumerate}
		\item Fecha
		\item Nombre
		\item Compañía
		\item Persona a quien visita
		\item Área que visita
		\item Asunto
		\item Hora de entrada y salida
		\item Firma
	\end{enumerate}
	\item Registrarse en la bitácora de entrada y salida de clientes el cual consta de los siguientes puntos:
	\begin{enumerate}
		\item Fecha
		\item Nombre
		\item Proveedor
		\item Cliente
		\item Identificación
		\item Hora de entrada y salida
		\item Folio de entrada y salida
	\end{enumerate}
	\item Si solo viene a las oficinas deberá esperar en recepción hasta ser atendido
	\item Si requiere de ingresar al área operativa deberá cumplir con el reglamento de ingreso y ser acompañado por la persona designada (Ver reglamento o política de entrada de áreas operativas).
	\item La persona ajena a la empresa será acompañado durante el total de su estancia dentro de la empresa
	\item Firmar la hoja de registro con la hora de salida
\end{itemize}

\subsubsection{Requisitos de ingreso para el personal}\index{Programa!Defensa Alimentaria, de!ingreso de personal}
\begin{itemize}
	\item Todo el personal deberá dejar sus pertenencias en el área de lockers (Esta área se encuentra fuera del área de operaciones) y ponerse su uniforme completo, este debe estar limpio.
	\item Antes de ingresar a las áreas operativas el personal debe colocar su lonche en el refrigerados para esta función (Esta área se encuentra fuera de las instalaciones de operación).
	\item El personal debe transitar por la aduana sanitaria y lavarse las manos (Ver reglamento de entrada de personal).
	\item En ningún momento se encuentra permitido el acceso dentro del almacén con mochilas o bolsas de mano y no se permite la entrada con ningún artículo personal dentro de las instalaciones.
	\item Al momento de la salida de planta el personal no puede llevar ningún artículo o producto al menos de que este sea autorizado por el jefe de área.
	\item El personal de vigilancia realiza inspección visual del personal y controla la salida del mismo
\end{itemize}

\subsubsection{Supervisión de empleados subcontratados}\index{Programa!Defensa Alimentaria, de!contratistas}
Toda persona que ingrese a la planta deberá cumplir ciertas reglas establecidas para su ingreso y permanencia en la misma.

\begin{itemize}
	\item Se registrara en la caseta de entrada y/o en oficinas administrativas.
	\item Esperará en recepción a la persona que lo va a atender
	\item Si requiere de ingresar al área operativa deberá cumplir con el reglamento de ingreso a planta.
	\item En todo momento deberá estar acompañado por la persona a la que visita.
	\item Para ingresar a cualquier otro departamento deberá tener la autorización y estar acompañado igual por personal de la planta
	\item Al final de visita la persona debe regresar a firmar la hoja de registro con la hora de salida en la caseta de vigilancia y/o recepción.
En el caso de proveedores que trabajan constantemente dentro de las áreas operativas se extiende un documento en el que la empresa confirma quien es el personal que presta sus servicios y que está autorizada para ingresar a la empresa.
\end{itemize}

\subsubsection{Contratación de personal}\index{Programa!Defensa Alimentaria, de!gestión de recursos humanos!contratación de personal}
\begin{itemize}
	\item Se cuenta con procedimiento para contratación de personal
	\item Se cuenta con el recopilado de información personal de la persona contratada.
	\item El personal se encuentra identificado por sistema, a través de número de nómina; como medio de reconocimiento de sus entradas y salidas a la planta.
	\item Al ser contratado el personal es ingresado a la base de datos de la empresa, se le designa un número de registro.
	\item Se informa al personal sobre el reglamento de trabajo de la empresa y BPD, así como los programas de calidad implementados.
\end{itemize}

\subsubsection{Notificación de baja del personal}\index{Programa!Defensa Alimentaria, de!gestión de recursos humanos!baja de personal}
\begin{itemize}
	\item De acuerdo a lo estableció en el plan de seguridad alimentaria lo importante es detectar alguna contaminación internacional en alguna parte del proceso; por lo que en el momento que se da de baja ya sea renuncia voluntaria o despido de algún trabajador, se da aviso de manera verbal a casetas de vigilancia donde se especifica el personal que ya no labora, con esto evitamos que personal que ya no labora en la empresa entre a las instalaciones sin la autorización requerida.
\end{itemize}

\subsubsection{Capacitación y entrenamiento del personal}\index{Programa!Defensa Alimentaria, de!gestión de recursos humanos!capacitación}
\begin{itemize}
	\item El personal de nuevo ingreso recibe capacitación de introducción, en el cual contiene temas de BPD, \gls{FD}, estos temas con la finalidad que el personal se encuentre entrenado en aspectos básicos del cómo actuar dentro de las áreas operativas y evitar acciones involuntarias de riesgo al producto.
	\item Se cuenta con programa de capacitación anual sobre temas que permitan el mejor desempeño y superación del personal que se verán reflejadas en los resultados, eficiencia y calidad en las actividades globales de las áreas operativas y por consecuencia en el producto que llegara al consumidor. Como resultados de este entrenamiento se busca el minimizar las acciones de peligro involuntarias que puedan poner el riego al producto.
	\item El personal recibe entrenamiento sobre \gls{FD} de manera anual como mínimo
\end{itemize}

\subsubsection{Sistema de monitoreo alarmas}\index{Programa!Defensa Alimentaria, de!sistemas de monitoreo}
En los ingresos se cuenta con sistema de alarma en puertas de acceso y sensores de movimiento en áreas de oficinas, las cuales son monitoreadas las \qty{24}{\hour} por el proveedor contratado.

\subsubsection{Vigilancia de actividad del personal}\index{Programa!Defensa Alimentaria, de!vigilancia}
\begin{itemize}
	\item Dentro de las áreas operativas se cuenta un sistema cerrado de cámaras (CIRCUITO CERRADO con grabación) que trabaja las 24 horas, estas se encuentra en diferentes puntos de las áreas operativas y ayudan a monitorear y controlar los movimientos dentro de las áreas operativas.
	\item La distribución de las cámaras es la siguiente:
	\item En el patio que abarca el área de andares y entrada principal a las áreas operativas
	\item En áreas de andares donde abarcan en total de los accesos principales de las áreas de refrigeración y congelación
	\item Personal autorizado se encarga del monitoreo de todas las cámaras para detectar alguna anormalidad.
	\item El área operativa cuenta con un jefe (Supervisor de almacén) que al mismo tiempo de coordinar las actividades de personal da aviso de cualquier conducta sospechosa.
	\item Anexado se encuentra un mapa de las áreas operativas planta donde se incluyen los puntos de localización de las cámaras.
	\item Las cámaras son verificadas de manera constante por el personal autorizado, las grabaciones quedan guardadas por 15 días. Las cuales pueden ser consultadas en el momento requerido.
\end{itemize}

\subsubsection{Control de llaves de ingreso al almacén}\index{Programa!Defensa Alimentaria, de!control de llaves}
\begin{itemize}
	\item Todos los accesos a las áreas operativas se encuentran cerrados con llave como medida de seguridad al final de las actividades operativas.
	\item Se tiene un llavero maestro al que solo tiene acceso el área de vigilancia y personal asignado para poder hacer uso de estas llaves.
	\item Las llaves están numeradas de acuerdo al flujo de entrada.
	\item En el llavero se encuentra una lista con el número y la descripción de la puerta a la que corresponde.
	\item Todas las llaves están restringidas para uso de personal no autorizado.
	\item Las llaves son entregadas únicamente al personal designado para cada área.
	\item Se cuenta con un registro de control de llaves. Cuando se realiza el préstamo de una llave para cuestiones operativas se debe de registrar.
	\item Las áreas restringidas se identifican en el procedimiento de áreas restringidas.
\end{itemize}

\subsubsection{Seguridad en proceso}\index{Programa!Defensa Alimentaria, de!seguridad en el proceso}
\begin{itemize}
	\item Se tiene implementado un sistema de Trazabilidad para rastrear producto recibido, almacenado, embarcado y distribuido.
	\item Se verifica la integridad de los productos y su empaque antes de su ingreso
	\item Se verifica la integridad de los productos y su empaque antes de su embarque.
	\item El inventario de producto terminado se lleva debidamente controlado y existen registros.
	\item Todos los productos se encuentran registrados en el sistema PROTHEUS.
\end{itemize}

\subsubsection{Áreas restringidas}\index{Programa!Defensa Alimentaria, de!áreas restringidas}
El personal responsable de áreas restringidas debe mantener estos accesos cerrados como son las áreas de:

\paragraph{Planta baja}\index{Programa!Defensa Alimentaria, de!áreas restringidas!planta baja}
\begin{enumerate}
	\item Acceso principal a patios de carga (Acceso restringido con barrera física portón y candado).
	\item Acceso Pasillo periferia de almacén (Acceso restringido con barrera física de puerta y llave).
	\item Puerta principal de acceso al almacén y Oficinas (Acceso restringido con barrera física puerta de seguridad y llave).
	\item Puerta principal de acceso a Oficina Supervisor (Acceso restringido con barrera física puerta de seguridad y llave).
	\item Medidor de Agua (Acceso restringido con barrera física)
	\item Registro de energía (Acceso restringido con barrera física de puerta y candado).
	\item Puertas de carga y descarga en el andén y rampas (Acceso restringido con barrera física puerta y barra de seguridad).
	\item Área artículos de limpieza (Acceso restringido con barrera física puerta y llave)
	\item Puerta de salida emergencia (Acceso restringido con barrera física puerta)
\end{enumerate}

\paragraph{Planta alta}\index{Programa!Defensa Alimentaria, de!áreas restringidas!planta alta}

\begin{enumerate}
	\item Área almacén de material de oficina y consumibles. (Acceso restringido con barrera física puerta y llave).
	\item Cuarto de controles de herramientas (Acceso restringido con barrera física puerta y llave).
	\item Cuarto de controles eléctricos (Acceso restringido con barrera física puerta y llave).
	\item Puerta principal de acceso a oficinas (Acceso restringido con barrera física puerta y llave).
\end{enumerate}

\subsubsection{Control de acceso para sistema de operación PROTHEUS}\index{Programa!Defensa Alimentaria, de!\gls{SGA}}

\begin{itemize}
	\item El acceso para el sistema de operación PROTHEUS para el control de producto recibido, almacenado, embarcado y distribuido se encuentran restringidos. Existe una lista de usuarios autorizados con una clave de acceso y son de los siguientes departamentos:
	\item Supervisor de almacén
	\item Mesa de Control
	\item Administrador General
	\item Facturación
	\item Gerencia de almacén
	\item El personal administrativo y de sistemas de cómputo son los responsables del control de autorizaciones al sistema.
\end{itemize}

\subsubsection{Entrada de Producto}\index{Programa!Defensa Alimentaria, de!ingreso de alimentos}

\begin{itemize}
	\item Se cuenta con inspección de transportes y productos (Ver formato de verificación de transporte), el cual consta de los siguientes puntos de verificación:
	\item Fecha
	\item Folio de verificación de registro
	\item Nombre del cliente
	\item Proveedor del cliente
	\item Temperatura de recibo
	\item Cantidad
	\item Número de lote
	\item Transporte
	\item Condiciones del transporte
	\item Libre de plaga
	\item Tarimas recibidas
	\item Observaciones
	\item Firmas
\end{itemize}

\subsubsection{Aéreas de productos de limpieza}\index{Programa!Defensa Alimentaria, de!áreas restringidas!almacen de productos de limpieza}

\begin{itemize}
	\item Los productos de limpieza se entregan por área específica y el pedido es controlado por el personal encargado del departamento de Mantenimiento de Edificio.
	\item Los productos de limpieza se mantienen como área de acceso restringido y bajo llave, solo el personal capacitado y autorizado para el manejo de productos de limpieza tiene acceso a esta área.
	\item El personal autorizado es el siguiente:
	\begin{itemize}
		\item Mantenimiento de Edificio.
		\item Mantenimiento de Equipos.
	\end{itemize}
\end{itemize}

\subsubsection{Salida de producto}\index{Programa!Defensa Alimentaria, de!salida de alimentos}

\begin{itemize}
	\item Se cuenta con un control de producto por lote durante su almacenamiento y de esta manera se controla el producto a embarcar, para fines de trazabilidad.
	\item Se cuenta con inspección de transportes y productos antes del embarque (Ver formato de verificación de transporte), el cual consta de los siguientes puntos de verificación:
	\item Fecha
	\item Folio de verificación de riesgo
	\item Nombre del cliente
	\item Proveedor del cliente
	\item Ruta
	\item Temperatura de embarque
	\item Pedido
	\item Maniobras de salida
	\item Picking
	\item Número de lote
	\item Transporte
	\item Condiciones del transporte
	\item Libre de plaga
	\item Observaciones
	\item Firmas
	\item Todas las entradas y salidas son registradas en el sistema interno de la empresa.
	\item El personal administrativo y de sistemas de cómputo son los responsables del control de autorización al sistema PROTHEUS.
	\item \gls{PPRO}
\end{itemize}

\subsection{Acciones correctivas}\index{Programa!Defensa Alimentaria, de!acciones correctivas}

\begin{itemize}
	\item Cuando se presenta una no conformidad en la realización del procedimiento, detectada por el supervisor o encargada del área se deberá reportar a su jefe directo
\end{itemize}

\subsection{Frecuencia}\index{Programa!Defensa Alimentaria, de!frecuencia}\index{Frecuencia!revisión del programa de \textit{food defense}}\index{Frecuencia!verificación del programa de \textit{food defense}}

\begin{itemize}
	\item[Revisión] Semestral;
	\item[Verificación] Semestral. 
\end{itemize}

\begin{changelog}[title=Registro de cambios,simple, sectioncmd=\subsection*,label=changelog-10.1]

	\begin{version}[v=\MayorVer.\MenorVer, date=2023--01, author=Pablo E. Alanis]
		\item Cambio de formato;
		\item Cambios en la serialización de versiones;
		\item Correcciones ortográficas y de estilo;
		\item Se agregó \cref{nota:MedicionDeTemperaturas}.
	\end{version}

	\begin{version}[v=1.7, date=2022--02, author=Alonso M.]
		\item cambio de fecha;
	\end{version}

	\shortversion{v=1.6, date=2021--02, changes=No hubo cambios}
\end{changelog}
\thispagestyle{formato-PI}
\renewcommand{\MayorVer}{2}
\renewcommand{\MenorVer}{0}
\renewcommand{\Codigo}{FD-1-MAN}
\renewcommand{\FechaPub}{2023--01}
%\renewcommand{\Edit}{2.1}
\renewcommand{\Titulo}{Equipo multidisciplinario de \textit{\glsname{FD}}}

\section{\Titulo} \label{ESP:EquipoFoodDefense}\index{Equipo multidisciplinario!\textit{\glsname{FD}}, de}

\subsection{Objetivo}
\begin{itemize}
	\item Crear el equipo multidisciplinario comprometido para asegurar el control de todos los puntos dentro de la cadena de recepción, almacenamiento y distribución de nuestros productos.
	\item Garantizar la seguridad (Inocuidad) de alimentos para consumo humano.
	\item Reducir, eliminar y evitar que ocurran riesgos significativos intencionales o no intencionales durante la operación.
\end{itemize}

\subsection{Alcance}
\begin{itemize}
	\item A todas las personas responsables para que se lleven a cabo las operaciones que se requieran para el buen cumplimiento de este procedimiento en RED DE FRIOS.
	\item Las pautas a seguir en este programa son aplicables a los productos, personal, instalaciones, visitantes, transporte y servicios externos contratados.
\end{itemize}

\subsection{Terminología y definiciones} %TODO

SEGURIDAD ALIMENTARIA: los alimentos por su naturaleza pueden ser blanco de un ataque. Por estar al alcance de mucha gente involucrada en su elaboración y distribución, de ahí que las Buenas Prácticas de Distribución y los sistemas de seguridad contribuyan a la prevención contra acciones deliberadas y mal intencionadas con los alimentos, dando lugar a los Programas de Bioseguridad Alimentaria.

\subsection{Documentos y/o normas relacionadas}

\begin{itemize}
	\item Programa de reuniones semestral.
	\item Procedimientos del Programa de Seguridad Alimentaria.
\end{itemize}

\subsection{Procedimiento}

La creación de un comité de \gls{FD} nos ayudara a proteger la seguridad del producto, apoyados en el personal de las diferentes áreas involucradas en la producción, calidad y almacenamiento y distribución.

\begin{itemize}
	\item El equipo debe ser multidisciplinario para diseñar un sistema global y comprometido preventivo no reactivo para dar la seguridad requerida durante la cadena de proceso y comercialización de producción y asegurar la inocuidad y seguridad de los alimentos.
	\item Este plan de Seguridad Alimentaria será desarrollado por un equipo multidisciplinario que puede constar de miembros de los siguientes departamentos.
	\begin{itemize}
		\item Director General
		\item Gerencia de Operaciones.
		\item Control de Inventarios y Coordinación del cliente.
		\item Mantenimiento a equipos.
		\item Mantenimiento a edificios.
		\item Supervisor de almacén.
		\item Recursos humanos.
		\item Sistemas.
	\end{itemize}
\end{itemize}

\subsubsection{Reuniones del equipo}

El equipo multidisciplinario se reunirá dos veces al año en el lugar previamente designado por el equipo para verificar como base los temas a continuación descritos.

\begin{itemize}
	\item Evaluar artículos de gravedad o riesgo visualizados durante la cadena de recepción, almacenamiento, embarque y distribución del producto.
	\item Comentar las tendencias en la industria de alimentos incluyendo los impactos sobre la seguridad alimentaria y bienestar público, estos temas pueden ser cambios o sugerencias aportados por cursos de capacitación sobre el tema.
	\item Revisar incidencias o acontecimientos significativos que se cataloguen como riesgo dentro de la cadena de recepción, almacenamiento, embarque y distribución del producto
	\item Revisar iniciativas reglamentarias o actualizaciones del sistema.
	\item Autorizar nuevos procedimientos de Seguridad Alimentaria que se adopten.
	\item Revisar los puntos de oportunidad del sistema de Seguridad Alimentaria para su correcta implementación.
	\item Actualización y elaboración de manuales del sistema si así se requiera.
	\item Puntos de oportunidad detectadas en auditorias de implementación del sistema de Seguridad Alimentaria.
	\item Análisis de observaciones y comentarios detectados por historial en el transcurso del año por parte de los integrantes del equipo.
	\item Programar y definir la capacitación del personal en materia de Seguridad Alimentaria.
\end{itemize}

\subsubsection{Registros}

\begin{itemize}
	\item Registrar en bitácora todos los puntos tratados en la reunión (Minutas).
	\item En la minuta se registran los puntos observados o áreas de oportunidad así mismo los responsables del cumplimiento de las acciones correctivas.
\end{itemize}

\subsection{Frecuencia}

Semestral

\begin{changelog}[simple, sectioncmd=\subsection*,label=changelog-10.2]

	\begin{version}[v=2.0, date=2023--01, author=Pablo E. Alanis]
		\item Cambio de formato;
		\item Cambios en la serialización de versiones;
		\item Correcciones ortográficas y de estilo;
		\item Se agregó \cref{nota:MedicionDeTemperaturas}.
	\end{version}

	\begin{version}[v=1.5, date=2022--02, author=Alonso M.]
		\item cambio de fecha;
	\end{version}

	\shortversion{v=1.4, date=2021--02, changes=No hubo cambios}
\end{changelog}
\thispagestyle{formato-PI}
\renewcommand{\MayorVer}{2}
\renewcommand{\MenorVer}{1}
\renewcommand{\Codigo}{BPD-4-CI/CE}
\renewcommand{\FechaPub}{2023--01}
\renewcommand{\Titulo}{Reglamento de ingreso a las instalaciones}

\section{\Titulo}\index{Reglamento!ingreso a las instalaciones, de}
%\section{Reglamento de ingreso a las instalaciones}

\subsection{Objetivos}
\begin{itemize}
	\item Garantizar la calidad e inocuidad de los productos alimenticios por medio del control del ingreso de personal y visitantes para evitar contaminación cruzada.
	\item Garantizar la eficiencia durante la recepción, almacenamiento y distribución de nuestros productos alimenticios.
\end{itemize}

\noindent
\subsection{Alcance}\index{Reglamento!ingreso a las instalaciones, de!alcance}
\begin{itemize}	
	\item Este documento se extiende a todo \gls{personal-interno} y \gls{cliente} de \gls{RDF}
\end{itemize}

\subsection{Términos y definiciones}

\begin{description}
	\defglo{personal-interno}
	\item[\gls{uniforme-completo}] \glsdesc{uniforme-completo}
	\item[\gls{cliente}] \glsdesc{cliente}
	\item[\gls{visitante}] \glsdesc{visitante}  
\end{description}

\subsection{Documentos y/o normas relacionados}
\begin{itemize}
	\item Programa de \textit{Food Defense;}
	\item Programa de \gls{HACCP};
	\item Programa de BPM;
	\item Programa de seguridad alimentaria;
	\item Ley Federal de Sanidad Animal (LFSA);
	\item Reglamento de Ley Federal de Sanidad Animal (RLFSA).
	\item Carta de conformidad con el reglamento interno
\end{itemize}

\subsection{Reglamento de ingreso al almacén}

\subsubsection{Disposiciones generales}\index{Reglamento!ingreso a las instalaciones, de!disposiciones generales}
\begin{itemize}
	\item Se deben de lavar las manos al entrar al área operativa con jabón y gel antibacteriano;\footnote{En caso de portar guantes, estos deberán estar limpios.}
	\item Traer su uniforme completo;
	\item Prohibido portar joyería;\footnote{Anillos, cadenas, aretes, pulseras, piercings en lugares visibles.}
	\item Prohibido fumar;
	\item Prohibido masticar chicle;
	\item Prohibido comer ó beber en las áreas operativas;
	\item Prohibido escupir;
	\item Usar ropa y/o uniforme limpio;
	\item Prohibido usar pantalones deshilachados o con pedrería;
	\item Prohibido entrar en shorts y camisas sin mangas;
	\item Usar cofias en áreas de dentro de Almacén;
	\item Usar cubrebocas en el Almacén;
	\item No portar artículos en bolsillos superiores a la cintura;
	\item Prohibido introducir artículos de vidrio;
	\item No golpear ni dañar los productos;
	\item Mantener limpias y en orden las áreas y las herramientas de trabajo;
	\item Prohibido el uso de zapatos abiertos;
	\item No pisar los productos ni las tarimas;
	\item No colocar en el piso materiales de contacto con producto \emph{i.e.} emplayes, cajas plástico, etc;
	\item No utilizar celulares personales dentro del almacén;\footnote{Si se tiene que utilizar el celular para una cuestión que involucre la operación, el personal que lo haya usado tendrá que lavarse y desinfectarse las manos nuevamente antes de tocar cualquier producto alimenticio.}
	\item Traer las uñas de las manos cortas;
	\item Prohibido usar barba larga;
	\item Si se tiene el cabello largo, éste debe de estar completamente cubierto por la cofia;
	\item Prohibido el uso de cámara fotográfica y de video de cualquier tipo dentro del almacén sin autorización previa.
\end{itemize}

\subsubsection{Ingreso de personal interno}\index{Reglamento!ingreso a las instalaciones, de!personal interno}
\begin{itemize}
	\item El personal contratado directamente de RDF debe de apegarse al reglamento interno;
	\item en caso de que se registre un incumplimiento, se deberá de registrar la insatisfacción y capacitar \emph{in situ} a la persona que incumplió con el reglamento.
\end{itemize}

\subsubsection{Ingreso de visitantes o personal contratado por clientes de RDF (PCPC)}\index{Reglamento!ingreso a las instalaciones, de!PCPC}
\label{PCPC-reglamento}

\begin{itemize}
	\item Los visitantes o personal contratado por clientes de RDF (PCPC) solo pueden ingresar a la planta de almacenamiento con previa autorización y acompañadas por personal autorizado;
	\item los visitantes y PCPC solo deberán tener acceso a las zonas que les son de incumbencia para su operación, por lo que no deben de permanecer sin estar acompañados, a la vista de personal de operaciones constantemente;
	\item los visitantes y PCPC se deben de apegar al reglamento descrito en el documento BPD-CI-1 y en el enumerado en la \cref{PCPC-reglamento};
	\item si los visitantes no portan el uniforme adecuado, no podrán ingresar a las áreas operativas y se les tratará de proporcionar uniforme adecuado.
	\begin{itemize}
		\item En caso de no contar con uniformes adicionales, se le notificará al cliente para determinar que acciones se tomarán.
	\end{itemize}
\end{itemize}

\subsection{Responsable(s) de la actividad}\index{Reglamento!ingreso a las instalaciones, de!responsables}
\begin{itemize}
	\item \emph{Ejecutado} por personal de operaciones y vigilancia;
	\item \emph{Monitoreado} por personal de operaciones;
	\item \emph{Verificado} por personal de gerencia y calidad.
\end{itemize}

\subsection{Acciones preventivas}

\begin{itemize}
	\item Se llevará a cabo una inspección por el personal de operaciones y en caso de ser observada una insatisfacción, se le pedirá al personal o visitante que incumplió con el reglamento cumplir con el requisito;
	\item El personal de calidad hará una inspección diaria de forma aleatoria y marcará en el formulario correspondiente si se cumple con el reglamento de ingreso de almacén.
	\item Si la desviación se repite frecuentemente se dará curso de capacitación al personal por parte de personal de calidad;
	\item Si después de haber capacitado al personal de almacén se siguen presentando desviaciones por causas injustificadas, será acreedor de una amonestación el personal interno que incumplió con el reglamento interno.
\end{itemize}

\subsection{Acciones correctivas}

\begin{itemize}
	\item En caso contrario la tarea se deberá volver a realizar como se indica en el procedimiento;
	\item Se prohíbe el acceso del personal re-incidente;
	\item En caso de no conformidad reportar en formato de acciones correctivas.
\end{itemize}

\begin{changelog}[simple, sectioncmd=\subsection*,label=changelog-1.4]
	\begin{version}[v=2.1, date=2023--01, author=Pablo E. Alanis]
			\item Cambio de formato;
			\item cambios en la serialización de versiones;
			\item cambio de Codificación;
			\item se contempló el acceso a PCPC.
	\end{version}

	\begin{version}[v=1.7, date=2022--05, author=Alonso M.]
		\item No hubo cambios.
	\end{version}

	\shortversion{v=1.6, date=2021--05, changes=No hubo cambios}
\end{changelog}

\subsection*{Anexos}
\begin{itemize}
	\item Ver \cref{carta.compromiso}
\end{itemize}
\newpage
\thispagestyle{formato-PI}
\renewcommand{\MayorVer}{2}
\renewcommand{\MenorVer}{0}
\renewcommand{\Codigo}{BPD-5-CI/CE}
\renewcommand{\FechaPub}{2023--01}
\renewcommand{\Titulo}{Reglamento de almacén}\index{Reglamento!almacén, de}
\section{\Titulo}

\begin{itemize}
	\item Mantener las manos limpias.
	\item Prohibido portar joyería (anillos, cadenas, aretes, pulseras, piercings en lugares visibles etc.)
	\item No se permite fumar
	\item No se permite masticar chicle
	\item Prohibido comer o beber en las áreas de Producto
	\item Prohibido escupir
	\item Usar ropa y/o uniforme limpio
	\item Prohibido usar pantalones deshilachados o con pedrería
	\item Prohibido entrar en shorts y camisas sin mangas
	\item Usar cofias en áreas de dentro de Almacén
	\item Usar cubrebocas en áreas de Almacén
	\item No portar artículos en bolsillos superiores a la cintura
	\item Prohibido introducir artículos de vidrio
	\item No golpear ni dañar los productos
	\item Mantener limpias y en orden las áreas y las herramientas de trabajo.
	\item Prohibido el uso de zapatos abiertos.
	\item No pisar los productos ni las tarimas
	\item No colocar en el piso material de contacto con producto (ejem.\ emplaye, cajas plástico, etc.)
	\item No utilizar celulares personales
	\item En caso de tener heridas expuestas o tener algúna enfermedad contagiosa, se debe de notificar previo a entrar al almacén.
\end{itemize}

\begin{changelog}[simple, sectioncmd=\subsection*,label=changelog-\thesection-\MayorVer.\MenorVer]
	\begin{version}[v=2.1, date=2023--01, author=Pablo E. Alanis]
			\item Cambio de formato;
			\item cambios en la serialización de versiones;
			\item cambio de Codificación.
	\end{version}

	\begin{version}[v=1.6, date=2022--05, author=Alonso M.]
		\item No hubo cambios.
	\end{version}

	\shortversion{v=1.5, date=2021--05, changes=No hubo cambios}
\end{changelog}
\newpage
\renewcommand{\MayorVer}{2}
\renewcommand{\MenorVer}{1}
\renewcommand{\Codigo}{BPD-6-CI}
\renewcommand{\FechaPub}{2023--01}
\renewcommand{\Titulo}{Sanción por no cumplir el reglamento de BPD}
\section{\Titulo}
\index{Sanciones!inclumplimiento de BPD}
%\section{Sanción por no cumplir el reglamento de BPD}

\subsection{Objetivo}

Establecer amonestación por incumplimiento del reglamento de buenas prácticas de distribución con el fin de hacer cumplir los reglamentos y procedimientos establecidos por la empresa.

\subsection{Alcance}

A todas las personas que incumplan los reglamentos y/o procedimientos establecidos por la empresa.

\subsection{Documentos y/o normas relacionados}

\begin{itemize}
	\item 1.3 - MA-OP-001 - Manual de calidad y buenas prácticas de distribución
	\item Tríptico de buenas prácticas de distribución.
\end{itemize}

\subsection{Procedimiento}

\subsubsection{Instrucciones}

\begin{itemize}
	\item Toda persona que ingrese a las áreas operativas deberá apegarse a los reglamentos de buenas prácticas de distribución establecidas por \gls{RDF} S.A de C.V.
	\item El no cumplimiento a los puntos establecidos será causa a una sanción administrativa.
	\item Se aplicara una amonestación por escrito a la persona que no cumpla con las buenas prácticas de manufactura.
	\item El no cumplimiento a las buenas prácticas de distribución, se le aplicara una amonestación dependiendo de la falta, estipulada en el capítulo IX de sanciones del Reglamento Interno de Trabajo.
	\item En caso de no pertenecer a la empresa (Visitante) se podrá revocar su permiso de ingreso.
	\item En caso de ser proveedor se podrá llegar el caso de cancelación de contratos.
\end{itemize}

\subsection{Responsable de la actividad}

\begin{itemize}
	\item \textbf{Ejecutado} por personal de RRHH;
	\item \textbf{Monitoreado} por personal de calidad;
	\item \textbf{Verificado} por personal de gerencia.
\end{itemize}

\subsection{Acciones preventivas}

\begin{itemize}
	\item Se llevara a cabo una inspección. Registrando el indicador y medida correspondiente
	\item Si la desviación se repite frecuentemente se dará curso de capacitación al personal.
	\item Si después de haber capacitado al personal de almacén se siguen presentando desviaciones por causas injustificadas, se tomaran acciones más enérgicas con el personal por incumplimiento con sus deberes.
\end{itemize}

\subsection{Acciones correctivas}

\begin{itemize}
	\item En caso contrario la tarea se deberá volver a realizar como se indica en el procedimiento.
	\item En caso de no conformidad reportar en formato de acciones correctivas.
\end{itemize}

\subsection{Frecuencia}
\index{Frecuencia!Sanciones}
\begin{itemize}
	\item Cuando el personal falte a la política de Buenas Prácticas de Distribución.
\end{itemize}

\begin{changelog}[simple, sectioncmd=\subsection*,label=changelog-1.6]
	\begin{version}[v=2.1, date=2023--01, author=Pablo E. Alanis]
			\item Cambio de formato;
			\item cambios en la serialización de versiones;
			\item cambio de Codificación;
			\item se contempló el acceso a PCPC.
	\end{version}

	\begin{version}[v=1.6, date=2022-05, author=Alonso M.]
		\item No hubo cambios.
	\end{version}

	\shortversion{v=1.5, date=2021-05, changes=No hubo cambios}
\end{changelog}

\subsection*{Anexos}

\begin{itemize}
	\item \cref{F-amonest}
\end{itemize}
\newpage
\thispagestyle{formato-PI}
\renewcommand{\MayorVer}{2}
\renewcommand{\MenorVer}{0}
\renewcommand{\TipoID}{PRO}
\renewcommand{\FechaPub}{2023--01}
\renewcommand{\TipoID}{PRO}
\renewcommand{\Titulo}{PEPS o VENCER}

\section{\Titulo}\index{Procedimiento!PEPS o VENCER}
\renewcommand{\Codigo}{\Prog--\thesection--\TipoID}
% \section{PEPS o VENCER}

\subsection{Objetivo}
Asegurar que los \glspl{alimento}, serán en sus completamente lotificados, identificados y verificados sus movimientos por medio de conteos programados de acuerdo a lo establecido en este procedimiento.

\subsection{Alcance}
Este documento será de utilidad para el personal del área de operaciones o el que se designe por la alta dirección como el encargado de identificar las tarimas de \glspl{alimento}

\subsection{Términos y definiciones}
\begin{description}
	\defglo{inventario-de-suministros-y-materiales} 
\end{description}
	
\subsection{Documentos y/o normas relacionadas}
\begin{itemize}
	\item Registros de recepción %TODO:Hacer Identificador
\end{itemize}

\subsection{Procedimiento}

\subsubsection{Instrucciones}
\paragraph{Primeras entradas, primeras salidas (PEPS) o VENCER}

\subparagraph{Ingreso de alimentos}
\begin{enumerate}
	\item Hacer la revisión correspondiente en cada recepción de producto (ver \cref{PRO:InspeccionDeAlimentosAlmacenados});
	\item revisar que se tengan todos los identificadores únicos del envío (número de lote, fechas, cantidades);
	\item regístralos en los formularios correspondientes para que se puedan ingresar en el \GLS{SGA}.
\end{enumerate}

\subparagraph{Etiquetado de alimentos}
\begin{enumerate}
	\item Se etiqueta por tarima y por lote cada producto;
	\item La etiqueta de \emph{Información de material} contiene los siguientes datos:
	\begin{itemize}
		\item Nombre de producto;
		\item Nombre del cliente;
		\item Fecha de recepción;
		\item Identificador de entrada;
		\item Lote;
		\item Cantidad.
	\end{itemize}
\end{enumerate}

\subparagraph{Consideraciones para el almacenaje de alimentos}
\begin{itemize}
	\item Cada \gls{alimento} con fecha de caducidad más tardía \emph{---Producto nuevo---}, se acomodará en la parte superior trasera de cada \textit{rack} y los \glspl{alimento} que ya estaban presentes se recorren a la parte inferior frontal.
	\item Al surtir los \glspl{alimento}, el personal de almacén deberá entregar los que tengan la \emph{fecha de recepción más antigua} o el producto que el cliente requiera para su operacion.
\end{itemize}

\begin{note}[Uso de \emph{PEPS}] \label{nota:PEPS}
\begin{itemize}
		\item El cliente puede decidir si se le surtirá \gls{alimento} en el esquema \emph{PEPS} o si requiere de un lote en específico;
		\item En caso de que el cliente prefiera que la salida de su producto sea por lote, y no por fecha de caducidad, \gls{RDF} no se responsabiliza si el cliente solicita lotes más nuevos aún cuando sigan existencias de lotes viejos o a punto de caducar almacenados en las instalaciones.\footnote{\gls{RDF} puede dar aviso a los clientes sobre mercancía a punto de caducar, sin embargo, \gls{RDF} no se responsabiliza de la destrucción de productos caducos, el cliente tiene que retirar \glspl{alimento} caducos.}
\end{itemize}
\end{note}

% \begin{lamp}[frametitle={Nota sobre uso de \emph{PEPS}}] \label{nota:PEPS}
% 	\begin{itemize}
% 		\item El cliente puede decidir si se le surtirá \gls{alimento} en el esquema \emph{PEPS} o si requiere de un lote en específico;
% 		\item En caso de que el cliente prefiera que la salida de su producto sea por lote, y no por fecha de caducidad, \gls{RDF} no se responsabiliza si el cliente solicita lotes más nuevos aún cuando sigan existencias de lotes viejos o a punto de caducar almacenados en las instalaciones.\footnote{\gls{RDF} puede dar aviso a los clientes sobre mercancía a punto de caducar, sin embargo, \gls{RDF} no se responsabiliza de la destrucción de productos caducos, el cliente tiene que retirar \glspl{alimento} caducos.}
% 	\end{itemize}
% \end{lamp}

\subparagraph{En el \glsfirst{SGA}}

\begin{itemize}
	\item En el \GLS{SGA} se cargan las entradas de producto en general al momento de su llegada.
	\item En este sistema administra la orden de salida de los mismos y permite dar de baja únicamente en el esquema de primeras entradas primeras salidas.
	\item El cliente también pudiera dar instrucciones del cual producto dar salida de forma prioritaria por estrategia comercial.
	\item En ambos casos el sistema de manera automática da de baja del sistema el producto que físicamente se da salida, actualizando las existencias.
\end{itemize}

\subsection{Responsables de la actividad}

\begin{itemize}
	\item \textbf{Ejecutado} por personal de operaciones;
	\item \textbf{Monitoreado} por personal de calidad;
	\item \textbf{Verificado} por personal de gerencia.
\end{itemize}

\subsection{Acciones preventivas}
\begin{itemize}
	\item Se llevará a cabo una inspección, registrando el indicador y medida correspondiente;
	\item Si la desviación se repite frecuentemente se dará curso de capacitación al personal de almacén y limpieza para que realicen eficientemente su trabajo.
	\item Si después de haber capacitado al personal de almacén se siguen presentando desviaciones por causas injustificadas, se tomaran acciones más enérgicas con el personal por incumplimiento con sus deberes.
\end{itemize}

\subsection{Acciones correctivas}

\begin{itemize}
	\item En caso contrario la tarea se deberá volver a realizar como se indica en el procedimiento.
	\item En caso de no conformidad reportar en formato de acciones correctivas (\RAC.
\end{itemize}

\subsection{Frecuencia}

\begin{itemize}
	\item Diaria.
\end{itemize}

\begin{changelog}[title=Registro de cambios,simple, sectioncmd=\subsection*,label=changelog-\thesection-\MayorVer.\MenorVer]

	\begin{version}[v=2.1, date=2023--01, author=Pablo E. Alanis]
		\item Cambio de formato;
		\item Cambios en la serialización de versiones;
	\end{version}

	\begin{version}[v=1.8, date=2022-02, author=Alonso M.]
		\item no se realizaron cambios de la revisión 06 a la 07.
	\end{version}

	\shortversion{v=1.7, date=2021--05, changes=No se realizaron modificaciones de la revisión 05 a la 06.}
\end{changelog}
\newpage
\thispagestyle{formato-PI}
\renewcommand{\MayorVer}{2}
\renewcommand{\MenorVer}{0}
\renewcommand{\Codigo}{PSA-2-PRO} 
\renewcommand{\FechaPub}{2023--01}
%\renewcommand{\Edit}{2.1}
\renewcommand{\Titulo}{Procedimiento de lotificación de productos para su almacenamiento}

\section{\Titulo}\index{Procedimiento!lotificación de productos}

% \section{Procedimiento de lotificación de productos para su almacenamiento}

\subsection{Objetivo}

\begin{itemize}
	\item \textbf{Asegurar} que las tarimas con producto a almacenar sean identificadas con la información requerida por el cliente.
\end{itemize}

\subsection{Alcance}

\begin{itemize}
	\item Este documento es aplicable a todas las tarimas con producto que almacenadas en RDF.
\end{itemize}

\subsection{Términos y definiciones}
%TODO definiciones

\subsection{Procedimiento}

\subsubsection{Fundamento}

Los establecimientos y equipos dedicados al proceso de alimentos para consumo humano deben mantener los registros y el control de los productos y materiales de empaque por medio de la notificación de los mismos, con el cual es posible la realización de un sistema de rastreo desde su producción hasta su distribución.

\subsection{Materiales}

\begin{itemize}
	\item Factura de entrada
\end{itemize}

\subsection{Instrucciones}

\begin{itemize}
	\item Cada vez que se reciba en el área de almacén, productos o materiales de empaque, se debe registrar el lote que el cliente designado a este producto.
	\item El personal de recepción debe de verificar que el lote de los registros coincidan con el lote marcado en las tarimas del producto;
	\item Estos se darán de alta en el sistema de la empresa con el lote asignado por el proveedor o cliente.
	\begin{itemize}
		\item Si el producto no tiene lote o este no coincide se dará aviso al Supervisor de almacén (el lote asignado será asignado con el número de entrada)
	\end{itemize}
	\item Las tarimas que contienen el producto son marcadas con \emph{etiqueta de recepción de producto,} en ella se coloca la fecha de recepción, nombre del cliente, nombre del producto, fecha de caducidad, orden de compra o lotificación, numero de entrada, cantidad por tarima.
\end{itemize}

\subsection{Responsables de la actividad}

\begin{itemize}
	\item \textbf{Ejecutado} por personal de operaciones;
	\item \textbf{Monitoreado} por personal de calidad;
	\item \textbf{Verificado} por personal de gerencia.
\end{itemize}

\subsection{Acciones preventivas}

\begin{itemize}
	\item Se llevará a cabo una inspección. Registrando el indicador y medida correspondiente;
	\item Si la desviación se repite frecuentemente se dará curso de capacitación al personal de almacén y limpieza para que realicen eficientemente su trabajo;
	\item Si después de haber capacitado al personal de almacén se siguen presentando desviaciones por causas injustificadas, se tomaran acciones más enérgicas con el personal por incumplimiento con sus deberes.
\end{itemize}

\subsection{Acciones correctivas}

\begin{itemize}
	\item En caso contrario la tarea se deberá volver a realizar como se indica en el procedimiento.
	\item En caso de no conformidad reportar en \RAC.
\end{itemize}

\subsection{Frecuencia}

\begin{itemize}
	\item Cada recepción de Producto
	\item Cada entrega de producto.
\end{itemize}

\begin{changelog}[simple, sectioncmd=\subsection*,label=changelog-2.2]
	
	\begin{version}[v=2.1, date=2023--01, author=Pablo E. Alanis]
			\item Cambio de formato;
			\item Cambios en la serialización de versiones;
	\end{version}

	\begin{version}[v=1.3, date=2022--05, author=Alonso M.]
		\item cambio de fecha;
		\item cambio de código.
	\end{version}

	\shortversion{v=1.2, date=2021--05, changes=No hubo cambios}
\end{changelog}
\newpage
\thispagestyle{formato-PI}
\renewcommand{\MayorVer}{2}
\renewcommand{\MenorVer}{0}
\renewcommand{\FechaPub}{2023--01}
\renewcommand{\TipoID}{PRO}
\renewcommand{\Titulo}{Inspección de producto durante su almacenamiento}

\section{\Titulo} \label{PRO:InspeccionDeAlimentosAlmacenados}\index{Inspección!productos durante su almacenamiento, de}
\renewcommand{\Codigo}{\Prog--\thesection--\TipoID}
\subsection{Objetivo}
Establecer un programa de inspección que asegure que las condiciones de almacenamiento se encuentren en condiciones óptimas como resultado de un adecudo mantenimiento de instalaciones, limpieza y temperatura apropiada para la conservación de productos alimenticios.

\subsection{Alcance}
\begin{itemize}
	\item Este documento se destina al área operativa, personal de mantenimiento y personal de aseguramiento de calidad.
\end{itemize}

\subsection{Términos y definiciones}
\begin{description}
	\defglo{alimento}
	\defglo{producto}
\end{description}

\subsection{Documentos y/o normas relacionados}
\begin{itemize}
	\item Programa de \glsfirst{BPD}.
\end{itemize}

\subsection{Procedimiento} \label{sec:pro:inspeccionDeProducto:almacenamiento}
\subsubsection{Materiales}
\begin{itemize}
	\item Linterna \textit{(Opcional).}
\end{itemize}

\subsection{Instrucciones}
\subsubsection{Temperatura de almacenamiento de alimentos refrigerados y congelados}
\begin{enumerate}
	\item Para asegurar que todo \gls{alimento} refrigerado o congelado se mantenga a la temperatura adecuada durante el almacenamiento, todos los productos se deben colocar directamente después de ser recibido en la cámara de refrigeración o congelación según la asignación, una vez que se descarguen.
	\item Las cámaras de refrigeración y de congelación están equipadas con un \emph{termómetro calibrado} ubicado en el \emph{rack,} en el cual se puede verificar la temperatura de cada cámara.
	\item Las lecturas de temperaturas deben ser documentadas \VecesTempManual veces durante el turno los días laborados.
	\item El Supervisor de mantenimiento es responsable de la verificación de temperaturas.
	\item Las puertas del andén de embarques deben permanecer cerradas en todo momento a menos que la unidad este siendo utilizada para carga y descarga de productos.
	\item Las puertas del andén deben de ser abiertas cuando la unidad de transporte ya se encuentre colocada en la rampa, esto con el fin de evitar fuga de temperatura y/o entrada de polvo o fauna nociva.
	\item Cualquier daño a las cámaras de refrigeración o congelación (techo, paredes, puertas, cortinas, etc.) debe ser reparadas en el menor tiempo posible.\footnote{El Supervisor de mantenimiento es responsable de garantizar que todo el equipo esté en buen estado y funcionando adecuadamente.}
\end{enumerate}

\subsubsection{Almacenaje y manejo de Productos Secos / Ambiente, Refrigerados y Congelados}
\begin{enumerate}
	\item Para asegurarse de que los productos almacenables a temperatura ambiente, refrigerados, y congelados son almacenados y protegidos de daños físicos, químicos, biológicos (contaminación), todos los productos deben ser almacenados \qty{\EspacioMinimoParedProducto}{\centi\meter} desde todas las paredes interiores\footnote{Si el diseño del almacén lo permite}, sobre atados limpios, \qty{\EspacioMinimoTechoProducto}{\centi\meter} de las superficies del techo, así como lejos de todas las luces superiores y nunca almacenados directamente en contacto con los pisos\footnote{todo \gls{alimento} almacenado en \gls{RDF} debe de estar sobre una tarima---no puede estar en contacto directo con el piso.}%TODO HACERLO INSTRUCCION 
	\item Todos los productos deben guardarse en su empaque secundario y etiquetados con las fechas de recepción necesarias.
	\item Todos los productos deben ser rotados de acuerdo en el principio de  \textit{“First-in-First-Out”} \textbf{(FIFO)}\footnote{En español, \gls{PEPS}} o primeras entradas primero a vencer (O según instrucciones de cada cliente).
	\item En caso de ser requerida, una \emph{área de retención} debe asignarse en el área de almacenamiento para segregar cualquier artículo recuperado, muestras de productos, producto dañado. \index{Área de retención}
	\item Todos los productos deben mantenerse libres de polvo y suciedad en todo momento.
	\item Cualquier o todos los productos químicos almacenados en la instalación deben estar completamente separados de todos los envases y productos alimenticios.
	\item Todos los productos químicos ubicados en la planta deben de tener un MSDS.
\end{enumerate}

\subsubsection{Control y mantenimiento de temperatura}
\begin{enumerate}
	\item Es imperativo que las cámaras en refrigerado y congelado en nuestras instalaciones se mantengan a la temperatura adecuada en todo momento. Esto incluye horas regulares, después de horas regulares, fines de semana y días festivos.
	\item Se tiene que hacer un seguimiento continuo de las temperaturas de las cámaras de almacenamiento, durante el horario normal las temperaturas son supervisadas y documentadas al día\footnote{Inicio de turno, medio turno y fin de turno} de lunes a sábado.
	\item El Supervisor de Mantenimiento es inmediatamente informado de cualquier irregularidad detectada en la temperatura,\footnote{Incluyendo horas después del trabajo, fines de semana y días festivos} de modo que las acciones correctivas se puedan tomar.
\end{enumerate}
\todo{hacerlo especificación}
\begin{scheme}
	\centering
	\missingfigure{Diagrama sobre inspección de alimentos}
	\caption[Inspección rutinaria de alimentos almacenados]{Inspección rutinaria de alimentos almacenados}
\end{scheme}

\begin{note}[Mediciones de temperaturas] \label{nota:MedicionDeTemperaturas} 
	Con motivos de controlar las temperaturas de las cámaras de almacenamiento, se cuentan con múltiples sistemas de medición para mantener una recursividad adecuada:
	\begin{itemize}
		\item El \emph{controlador de las unidades de refrigeración} es la principal fuente de información sobre las temperaturas de las cámaras. La información que el controlador recaba se basa en los termopares individuales de cada cámara, estos están debidamente considerados en el \emph{plan de metrología} y son calibrados anualmente;
		\item Semanalmente se colocan \emph{termograficadores} en las cámaras de almacenamiento y se verifican las temperaturas por el departamento de aseguramiento de calidad.
		\item Diariamente se comprueban las temperaturas de las cámaras de forma manual con termómetros infrarrojos por parte de:
		\begin{itemize}
			\item aseguramiento de calidad;
			\item mantenimiento.
		\end{itemize}
	\end{itemize}
\end{note}

\todo{hacerlo instruccion}

\subsubsection{Plan de contingencia en caso de avería de equipos de refrigeración}\index{Plan de contingencia!en caso de avería de equipos de refrigeración}
Con fin de garantizar que todos los productos refrigerados y congelados se encuentren protegidos contra cualquier peligro microbiológico asociado con el incremento de la temperatura, las cámaras con refrigeración o congelación que presenten una falla y que no puedan ser reparadas en menos en el caso de refrigeración de \qty{\TiempoAveriaRefri}{\hour} y en congelación de \qty{\TiempoAveriaConge}{\hour},\footnote{Siempre y cuando las puertas de cámara se encuentren cerradas y el producto mantenga la temperatura establecida} el producto se \emph{debe} trasladar oportunamente a otra cámara previamente acondicionada a las especificaciones del producto.

\subsection{Responsables de la actividad}

\begin{itemize}
	\item \textbf{Ejecutado} por personal de operaciones y mantenimiento;
	\item \textbf{Monitoreado} por personal de calidad;
	\item \textbf{Verificado} por personal de gerencia.
\end{itemize}

\subsection{Acciones preventivas}\todo{actualizar}
% \AcPGen

\subsection{Acciones correctivas}
\begin{itemize}
	\item En caso contrario la tarea se deberá volver a realizar como se indica en el procedimiento.
	\item En caso de no conformidad reportar en formato de acciones correctivas \RAC.
\end{itemize}

\subsection{Frecuencia}

\begin{itemize}
	\item \textbf{Verificación manual de temperaturas:}
		\begin{enumerate}
			\item \emph{Dpto.\ de mantenimiento:} \VecesTempManual veces al día;
			\item \emph{Dpto.\ de aseguramiento de calidad:} \VecesTempManual veces al día;
		\end{enumerate}
\end{itemize}

\begin{changelog}[title=Registro de cambios,simple, sectioncmd=\subsection*,label=changelog-\thesection-\MayorVer.\MenorVer]

	\begin{version}[v=2.1, date=2023--01, author=Pablo E. Alanis]
		\item Cambio de formato;
		\item Cambios en la serialización de versiones;
		\item Correcciones ortográficas y de estilo;
		\item Se agregó \cref{nota:MedicionDeTemperaturas}.
	\end{version}

	\begin{version}[v=1.8, date=2022--05, author=Alonso M.]
		\item cambio de fecha;
	\end{version}

	\shortversion{v=1.7, date=2021--05, changes=No hubo cambios}
\end{changelog}
\newpage
\thispagestyle{formato-PI}
\renewcommand{\MayorVer}{2}
\renewcommand{\MenorVer}{0}
 %TODO
\renewcommand{\FechaPub}{2023--01}
\renewcommand{\TipoID}{L}
\renewcommand{\Titulo}{Lista de contactos de emergencia}

\section{\Titulo}\index{Lista!contactos de emergencia, de}
\renewcommand{\Codigo}{\Prog--\thesection--\TipoID}

\subsection{Objetivo}
\begin{itemize}
	\item Crear el equipo para actuar en casos de crisis por corte de energía eléctrica, desastres naturales o siniestros.
	\item Generar la coordinación del plan de contingencia a la etapa de crisis
	\item Garantizar la seguridad e inocuidad de alimentos para consumo humano.
\end{itemize}

\subsection{Alcance}
A todas las personas responsables para que se lleven a cabo las operaciones que se requieran para el buen cumplimiento de este procedimiento en \gls{RDF}

\subsection{Términos y definiciones}

\begin{description}
	\item[equipo de emergencia] el grupo de personas responsables de desarrollar, implementar y verificar el cumplimiento del programa de emergencia.
\end{description}

\subsection{Documentos y/o normas relacionadas}
\begin{itemize}
	\item Procedimiento en caso de corte de energía o desastre natural o siniestro. \cref{sec:AccEnCasoDeEmergencia}
\end{itemize}

\subsection{Procedimiento}
\subsubsection{Equipo de Emergencias:}
\paragraph{Coordinador 1}
Gerardo Ruiz, Gerente de Operaciones.
Teléfono: +52 (81) 8351--0512, +52 (81) 8029--5609
Responsable de coordinar los esfuerzos según la emergencia y restablecer la operación, garantizando al mismo tiempo las políticas más estrictas de seguridad y protocolos.

\paragraph{Coordinador 2}
Luis Escareño, Jefe de Almacén
Teléfono: +52 (81) 8351--0512, +52 (81) 1390--7099
Responsable de asegurar el entorno de almacén y asegurarse de que el personal de almacén salga del edificio inmediatamente y se congreguen en el exterior en el punto de reunión en caso de emergencia.

\paragraph{Coordinador 3}
Salvador Verdín, Jefe de Mantenimiento
Teléfono: +52 (81) 8351--0512, +52 (81) 1706--1644.
Responsable del personal de oficinas salga del edificio inmediatamente y se congreguen en el exterior en el punto de reunión en caso de emergencia.

\paragraph{Coordinador 4}
Mario González, Gerente de Mantenimiento
Teléfono: +52 (81) 1292--7272, +52 (81) 1607--6035
Responsable de asegurar las condiciones internas de almacén y equipos, asegurar que el personal de almacén salga del edificio inmediatamente y se congreguen en el exterior en el punto de reunión en caso de emergencia.

\paragraph{Vigilancia}

Lorenzo Hernández, Vigilancia
Teléfono: +52 (81) 2353--2346

\paragraph{Números de Contacto de Emergencia}
\begin{itemize}
	\item Además del área de vigilancia, cada instalación debe de tener los números de emergencia de las autoridades locales, estatales y federales.
	\item Responsable de verificar de primera instancia alguna anormalidad vista desde el punto externo del almacén y dar aviso al personal correspondiente según el hallazgo detectado.
\end{itemize}

\begin{changelog}[simple, sectioncmd=\subsection*,label=changelog-\thesection-\MayorVer.\MenorVer2]
	\begin{version}[v=\MayorVer.\MenorVer, date=2023--01, author=Pablo E. Alanis]
		\item Cambios de formato;
		\item cambio de código;
		\item se añadieron recuadros con los contactos y contratos de la \gls{CFE}.
	\end{version}

	\begin{version}[v=1.9, date=2022--01, author=Agustín M.]
		\item febrero 2022 se realizó un acomodo a la información y se agregaron los No. De contrato con CFE de la revisión 07 a la 08.
	\end{version}
\end{changelog}

\newpage
\thispagestyle{formato-PI}
\renewcommand{\MayorVer}{1}
\renewcommand{\MenorVer}{0}
\renewcommand{\Codigo}{FD-ToC}
\renewcommand{\FechaPub}{2023--01}
%\renewcommand{\Edit}{01}
\renewcommand{\Titulo}{Programa de \textit{Food Defense} --- Glosario}

\printglossary[type=acronym]
\printglossary
\thispagestyle{formato-PI}
\renewcommand{\MayorVer}{1}
\renewcommand{\MenorVer}{0}
\renewcommand{\Codigo}{FD-IND}
\renewcommand{\FechaPub}{2023--01}
%\renewcommand{\Edit}{01}
\renewcommand{\Titulo}{Programa de \textit{Food Defense} --- Indice alfabético}

\printindex
\end{document}